\section{Operadores de masa en HMT--MD: escala absoluta, atlas de rutas y realización por sectores}
\label{sec:tabla-completa-masas-hmt-sm-revfinal}

La masa aparece en Mecánica Dimensional como la realización de dos ramas
matemáticas que confluyen sin confundirse. La primera conserva la genealogía del
estado: una extensión superviviente determina una ruta observable; la ruta escribe
un registro cronológico dodecafásico; el registro produce una firma entera, y la firma
alimenta el carácter y la jerarquía armónica. La segunda construye la recta
dimensional: el lector determinantal fija la escala de acción, el reloj interno la
transporta a energía y la carta cinemática la lleva a masa. El operador de especie
actúa entonces por homotecias positivas sobre esa recta. La coordenada experimental
se abre solamente después de congelar ambas ramas.

El dominio interno completo posee cuatro censos distintos:
\begin{equation}
 81=25_{\ell^-}+20_{\nu}+20_d+16_u,\qquad
 56=41\!\times\!1+10\!\times\!2+2\!\times\!3+1\!\times\!4+2\!\times\!5,
 \label{eq:censos-81-56-masas-rectoras}
\end{equation}
\begin{equation}
 13\ \text{multisecciones},\qquad
 324=81\times4\ \text{rutas orientadas}.
 \label{eq:censos-13-324-masas-rectoras}
\end{equation}
Las realizaciones físicas comparadas más adelante constituyen una
clasificación de realizaciones; no son otro nombre para las 81 celdas, las 56
familias, las 13 multisecciones o las 324 historias orientadas. Sobre este dominio
los dos lectores de ruta producen catorce perfiles \(K_D\), nueve perfiles
\(K_\Omega\), cuarenta pares conjuntos y dieciocho coincidencias. El sector nulo
permanece separado del carácter positivo de masa.

\subsection{De la incompatibilidad APP a la firma de masa}

La diferencia primordial entre las hojas de suma y producto aparece en la
descomposición espectral del bloque central de APP:
\begin{equation}
 B_c=7I+2R=9P_++5P_-,\qquad
 \operatorname{Spec}(B_c)=\{9,5\},\qquad
 \Delta_{\rm APP}=9-5=4.
 \label{eq:salto-app-cuatro-ley-masas}
\end{equation}
Los enteros que reaparecen después conservan fuentes distintos. El coeficiente
local de intercambio es \(2\) en \(B_c=7I+2R\); la compresión modular publica
\(51-45=6\), y su despliegue sobre nueve bloques produce \(9\cdot6=54\). Estas
igualdades pertenecen a niveles tipados de una misma genealogía y no constituyen
una cadena de definiciones del salto \(4\). En la sección electrónica, \(54\) es
además \(D_1(\Gamma_e)\) y la periodicidad cinemática de orden \(54\) del
endomorfismo temporal CT108: mide una componente de la historia ya transportada
por TPK, no la separación espectral primordial.

Para cada ruta enriquecida \(\Gamma\), el registro dodecafásico determina
\begin{equation}
 \sigma(\Gamma)
 =(n_A,s_C,k_{\Delta_4},b_{120},b_\zeta)\in\mathbb Z^5,
 \qquad \zeta_{270}=135\Delta_4^2.
 \label{eq:firma-z5-masas-rectoras}
\end{equation}
El carácter central y su refinamiento armónico son
\begin{align}
 \log\chi_0(\sigma)
 &=n_A A_{\rm rad}+\frac{s_C}{6}C^*_{\rm rad}
   +k_{\Delta_4}\Delta_4+b_{120}s_{120}+b_\zeta\zeta_{270},
 \label{eq:caracter-central-masas-rectoras}\\
 \log\chi_{\rm full}(\sigma,\eta)
 &=n_A A_{\rm rad}+\frac{s_C}{6}C^*_{\rm rad}
   +k_{\Delta_4}\Delta_4+b_\zeta\zeta_{270}
   +(b_{120}+\eta_{120})s_{120}
   +\sum_{\substack{h\in\mathfrak S\\h\ne120}}\eta_hs_h.
 \label{eq:caracter-completo-masas-rectoras}
\end{align}
La división \(s_C/6\) procede de la integralidad dodecafásica. El par etiquetado
\((b_{120},\eta_{120})\) conserva su genealogía incluso cuando su evaluación se
agrupa en un único coeficiente.

Los momentos que intervienen en
\eqref{eq:caracter-completo-masas-rectoras} pertenecen a una jerarquía infinita:
\begin{equation}
 q_\pm=\exp\!\left[-\frac{\pi}{180}(A\pm C^*)\right],
 \qquad s_h=q_-^h-q_+^h,\qquad c_h=q_-^h+q_+^h,
 \label{eq:momentos-torres-masas-rectoras}
\end{equation}
\begin{equation}
 \mathfrak S=\langle90,120\rangle
 =\{90,120\}\cup\{30r:r\ge6\}.
 \label{eq:semigrupo-global-torres-masas-rectoras}
\end{equation}
Tanto \(s_n\) como \(c_n\) satisfacen
\begin{equation}
 X_{n+2}=c_1X_{n+1}-(q_-q_+)X_n.
 \label{eq:recurrencia-torres-masas-rectoras}
\end{equation}
Las alturas \(90,120,180,210,240,270,360,420\) son testigos genealógicos
particularmente visibles; el dominio de la ley es el semigrupo completo. En
particular, \(360\) conserva dos historias antes del cociente, cuatro pasos de
\(90\) y tres de \(120\), aun cuando ambas produzcan el mismo momento escalar.

\subsection{La sección absoluta de masa y la década \texorpdfstring{\(10^{-34}\)}{10⁻³⁴}}

El bloque local de Hadamard
\[
 H_4=
 \begin{pmatrix}
  1&1&1&1\\
  1&1&-1&-1\\
  1&-1&1&-1\\
  1&-1&-1&1
 \end{pmatrix},
 \qquad
 \operatorname{SNF}(H_4)=\operatorname{diag}(1,2,2,4),
\]
tiene cokernel de orden \(16\). Su lector determinantal produce
\begin{equation}
 \Sigma_H=\varphi\alpha_{\rm HMT}^{16},
 \qquad
 10^{-34}<\Sigma_H<10^{-33},
 \qquad
 \left\lfloor\log_{10}\Sigma_H\right\rfloor=-34.
 \label{eq:orden-menos-34-tabla-masas}
\end{equation}
La conversión a \(\mathrm{J\,s}\) es una carta posterior: la potencia
\(10^{-34}\) nace en la estructura discreta y fija la escala de una única recta
de acción. Sus secciones anterior y posterior al retorno son
\begin{equation}
 \hbar_{\rm pre}=H_5\,10^{-34}\mathcal U_S,\qquad
 \hbar_{\rm ret}=\eta_{\rm ret}\,10^{-34}\mathcal U_S,
 \qquad
 R_{\rm act}=\frac{H_5}{\eta_{\rm ret}}>0.
 \label{eq:accion-absoluta-y-razon-tabla-masas}
\end{equation}
El desdoblamiento posee una carta aditiva,
\(\delta_{2w}=H_5-\eta_{\rm ret}\), y la carta multiplicativa adimensional
\(R_{\rm act}\). Declarada una duración interna \(t_0\), la periodicidad de
108 iteraciones del endomorfismo temporal CT108 construye
\begin{align}
 \omega_0&=\frac{2\pi}{108t_0},
 &\varepsilon_{\rm clk}&=\hbar_{\rm ret}\omega_0,
 \label{eq:cuanto-reloj-masa-rectora}\\
 \mathbf m_\star^{\rm pre}
 &=\left(H_5\,10^{-34}\mathcal U_S\right)
   \frac{2\pi}{108t_0}c_{\rm int}^{-2},
 &\mathbf m_\star^{\rm ret}
 &=\left(\eta_{\rm ret}\,10^{-34}\mathcal U_S\right)
   \frac{2\pi}{108t_0}c_{\rm int}^{-2}\in\mathcal L_M.
 \label{eq:seccion-reloj-masa-rectora}
\end{align}
Ambas secciones satisfacen
\(\mathbf m_\star^{\rm pre}/\mathbf m_\star^{\rm ret}=R_{\rm act}\).
La década \(10^{-34}\), el reloj común y \(c_{\rm int}^{-2}\) fijan la escala
absoluta. En una razón de masas esa sección común cancela exactamente; el
refinamiento relativo pertenece a la ruta y no a una segunda introducción de la
potencia decimal.

\subsection{Lectores bidireccionales de trayectoria y ley operatorial sobre las \(13\) multisecciones}

Los lectores que actúan sobre una ruta \(\Gamma\) son
\begin{equation}
 K_D(\Gamma)
 =\left(D_{27}+D_{36},\,D_1,\,\frac{D_4-D_3}{2}\right),
 \qquad
 K_\Omega(\Gamma)=(\Omega,54,P_6).
 \label{eq:lectores-kd-komega-masas-rectoras}
\end{equation}
Cada lector induce sobre una multisección \(F\) un operador diagonal
\(\widehat K_F^r\), \(r\in\{D,\Omega\}\). El índice \(s\) designa una sección o
bloque físico dentro de \(F\), y \(\widehat B_{F,s}^r\) es la restricción
autoadjunta de \(\widehat K_F^r\) a ese bloque. Si
\(\widehat{\mathcal A}^{(0)}_{F,s}\) es la restricción del operador que reúne la
firma, el carácter completo y la jerarquía global de torres, la ley general toma
la forma
\begin{equation}
 \boxed{
 \widehat{\mathbf M}^{\,\beta,r}_{F,s}
 =\mathbf m_\star^{\rm ret}\otimes
  R_{\rm act}^{\widehat B_{F,s}^r/2}
  \widehat{\mathcal A}^{(0)}_{F,s}
  R_{\rm act}^{\widehat B_{F,s}^r/2},
 \qquad r\in\{D,\Omega\}.}
 \label{eq:ley-operatorial-general-masas-rectoras}
\end{equation}
La escritura simétrica preserva positividad sin imponer la conmutación de
\(\widehat{\mathcal A}^{(0)}_{F,s}\) y \(\widehat B_{F,s}^r\). En la base
simultáneamente diagonal de los certificados vigentes coincide con el producto
\(\mathbf m_\star^{\rm ret}\otimes\widehat{\mathcal A}^{(0)}_F
\exp((\log R_{\rm act})\widehat K_F^r)\).
Esta ecuación gobierna las 81 celdas, las 56 familias, las 13 multisecciones y las
324 rutas orientadas. Fuera de una fibra de rango uno, la salida primaria es el
operador completo y su espectro. La realización escalar usa el morfismo cerrado
\(\mathcal S_{\rm phys}:(\sigma_{\rm ud},g,\chi_{\rm fina},[\gamma]_{\rm col})
\mapsto\gamma^{\rm surv}_{P,\mathfrak M}\) de
\cref{eq:realizacion-fisica-sabor-ruta-cerrada-rev2}; esa selección es
independiente de masas objetivo y antecede a cualquier comparación metrológica.

\subsection{Reducción de rango \(1\) en la construcción del electrón central}

La sección electrónica es el único punto fijo central del atlas:
\begin{equation}
 (r,c)=(5,5),\qquad
 \Gamma_e=(5,5,++),\qquad
 \tau_e=\mathtt{+++---+++---}.
 \label{eq:ruta-electronica-central-masas-rectoras}
\end{equation}
Sobre ella coinciden ambos lectores,
\begin{equation}
 K_D(\Gamma_e)=K_\Omega(\Gamma_e)=(80,54,6),\qquad
 \kappa_e=80-54\alpha_{\rm HMT}+6\Delta_4.
 \label{eq:kappa-electron-rector}
\end{equation}
La reducción de rango uno de
\eqref{eq:ley-operatorial-general-masas-rectoras} entrega
\begin{equation}
 \mathbf m_e^\beta
 =\mathbf m_\star^{\rm ret}\,
   \mu_e^{(0)}R_{\rm act}^{\,80-54\alpha_{\rm HMT}+6\Delta_4},
 \label{eq:electron-rank1-ley-masas}
\end{equation}
cuya coordenada en la carta comparativa es
\begin{equation}
 \boxed{m_e^\beta
 =0.5109989506900008086082571648\ldots\ \mathrm{MeV}/c^2.}
 \label{eq:masa-electron-tabla-completa-revfinal}
\end{equation}
El valor basal
\(0.5109989224880660725\ldots\ \mathrm{MeV}/c^2\) conserva su función como
etapa intermedia; la salida beta de
\eqref{eq:masa-electron-tabla-completa-revfinal} gobierna la comparación final.

\subsection{Clasificación física y comparación cuantitativa por sectores}

La exposición se organiza en cuatro estratos coordinados. El primero publica el
resultado interno y su codominio; el segundo declara el morfismo de realización
física; el tercero muestra la mejor coordenada numérica ya disponible y conserva
sus hipótesis y su dominio; el cuarto introduce el dato externo como contraste posterior.
Cada columna conserva así su tipo matemático: salida HMT, morfismo de
realización, coordenada calculada y contraste externo. La tabla publica
simultáneamente las evaluaciones disponibles y los mapas de construcción que las
producen.

El electrón publica su clausura beta escalar. El fotón y los gluones publican
secciones nulas exactas en sus cartas de masa, junto con extensiones gauge no
triviales. Los otros leptones cargados y los seis quarks publican simultáneamente
sus rutas nominales, caracteres escalares, operadores diagonales y espectros sobre
sus fibras; los neutrinos publican el operador neutro y su mezcla. \(W\), \(Z\) y
\(H\) poseen rutas CT--108 individualizadas, huellas completas y caracteres
\((n_A,s_C)=(94,-1),(95,-1),(97,9)\). El par \(K_D/K_\Omega\) pertenece al
refinamiento de las trece multisecciones del atlas y no es una premisa de la
realización bosónica por ruta nominal. Compuestos, sectores topológicos y estados
virtuales mantienen respectivamente sus codominios de ligadura, fusión, trenzado
y obstrucción.

\Needspace{15\baselineskip}
\subsubsection{Correspondencia tipada entre objetos internos y codominios físicos}

